\documentclass[9pt]{developercv} % Default font size, values from 8-12pt are recommended

% es-noquoting stops babel-spanish making " active, which would collide with the
% template's markup. Without babel this document was hyphenated with ENGLISH patterns.
\usepackage[spanish,es-noquoting]{babel}

% Personal data, shared by resume.tex (EN) and curriculum.tex (ES).
%
% Everything here is language-independent, so a changed phone number, handle or
% keyword is edited ONCE instead of once per language -- the two documents used to
% carry byte-identical copies of all of it, which is exactly the kind of duplication
% that goes stale in one file only.
%
% The link macros wrap the whole \href rather than just the URL on purpose:
% hyperref reads its URL argument under modified catcodes, so a macro placed inside
% \href{...} is not reliably expanded. Written this way, \href always sees literal
% text. The '\_' escapes are likewise required at definition time -- the macro body
% is tokenised when it is defined, before hyperref can adjust anything.

\newcommand{\cvfirstname}{Alin}
\newcommand{\cvlastname}{Trandafir}

\newcommand{\cvphone}{+34 653 751 742}
\newcommand{\cvemail}{\href{mailto:alin\_trandafir@icloud.com}{alin\_trandafir@icloud.com}}
\newcommand{\cvweb}{\href{https://end.works}{end.works}}
\newcommand{\cvgithub}{\href{https://github.com/ender-null}{github.com/ender-null}}
\newcommand{\cvlinkedin}{\href{https://linkedin.com/in/alin-trandafir}{linkedin.com/in/alin-trandafir}}

% Technology keywords for PDF /Keywords metadata. Proper nouns only, so this list is
% identical in both languages; each document appends its own translated terms.
% Legitimate metadata, not hidden body text -- do NOT put invisible keywords in the
% page content.
\newcommand{\cvkeywords}{%
	TypeScript, JavaScript, React, React Native, Angular, Next.js, Node.js,
	Expo, EAS, NestJS, Redux, RxJS, Firebase, Docker, Cypress, Jenkins, GitLab CI, CI/CD,
	Java, Symfony, Stripe, MapBox, i18next, i18n, l10n,
	Google Cloud, GCP, REST, SaaS, MongoDB, MySQL, SQL, NoSQL, Firestore,
	Keychain, Keystore, MFA, Agile, Scrum, CEO, Frontend, Zaragoza,
	Claude, Claude Code, LLM, MCP, AI, GenAI%
}
 % Nombre, datos de contacto y lista compartida de palabras clave

% Metadata. pdflang=es-ES also gives parsers the signal to load a Spanish month table
% -- "Dic" is the one month abbreviation here that an English table does not resolve.
\hypersetup{
	pdftitle={\cvfirstname\ \cvlastname\ -- Desarrollador Full-Stack Senior -- CV},
	pdfauthor={\cvfirstname\ \cvlastname},
	pdfsubject={Desarrollador Full-Stack Senior -- TypeScript, React, Angular, React Native},
	pdfkeywords={\cvkeywords, internacionalizacion, cifrado, autenticacion,
		Cofundador, Fundador, Desarrollador Full-Stack, Movil, Espana, Remoto},
	pdflang={es-ES},
}

% "Nov 2024 -- Actualidad" needs 92.8pt; the 0.20 default gives 92.2pt and a \parbox
% overflows silently, running the date into the job title.
\setlength{\dateColWidth}{0.23\textwidth}

\begin{document}

% Bloque de nombre, puesto y datos de contacto. La maquetación está en developercv.cls
% y los valores en identity.tex; aquí solo van las tres cadenas traducidas.
% El puesto se parte en dos líneas a mano: no cabe en una sola en la columna izquierda.
\cvheader{Desarrollador Full-Stack\\Senior}{Zaragoza, España}{Teléfono}

\cvsect{Perfil Profesional}

% "Perfil Profesional" rather than "¿Quién soy?": ATS segment a document by matching
% section headers against a dictionary, and "¿Quién soy?" matches nothing in any of them.

\textbf{Desarrollador full-stack senior} con \textasciitilde{}9 años de experiencia en \textbf{TypeScript}, \textbf{React}, \textbf{Angular} y \textbf{React Native}, e \textbf{inglés C1} certificado (Oxford Test of English, 2025). He construido aplicaciones en producción y plataformas web escalables en grandes empresas, y desarrollo y opero las mías propias: \textbf{Zierzo} y \textbf{LoveLust} están publicadas en App Store y Google Play a través de \textbf{The Lovelust Company, SL}, la empresa que cofundé y dirijo como CEO, con LoveLust alcanzando los \textbf{16.000 usuarios activos mensuales}.
\\
Desarrollo \textbf{con IA integrada en mi flujo de trabajo}: uso a diario \textbf{Claude}, LLMs y \textbf{MCP} para desarrollar, automatizar pruebas y generar assets y contenido.
\\
Me siento cómodo \textbf{liderando la dirección técnica} en todo el stack, y me apasiona crear software que \textbf{facilite la vida de las personas}.

\cvsect{Competencias Técnicas}

\textbf{Lenguajes:} TypeScript, JavaScript (ES6+), Java, PHP, HTML5, CSS3\\
\textbf{Frontend:} React, Angular, Next.js, Redux, RxJS, gestión de estado, librerías de componentes, diseño responsive\\
\textbf{Móvil:} React Native, Expo, EAS, Cordova, Keychain de iOS / Keystore de Android, publicación en App Store y Google Play\\
\textbf{Backend:} Node.js, NestJS, APIs REST, Symfony\\
\textbf{Bases de datos:} MongoDB, MySQL, SQL, Firebase / Firestore\\
\textbf{Cloud y DevOps:} Google Cloud Platform (GCP), Docker, CI/CD, GitLab CI, Jenkins, Git\\
\textbf{Desarrollo con IA:} Claude (Claude Code), LLMs, MCP (Model Context Protocol), automatización de pruebas con IA, generación de assets y contenido\\
\textbf{Testing y Calidad:} React Testing Library, Playwright, Cypress, tests unitarios, tests end-to-end, cobertura de tests, code review\\
\textbf{Metodologías:} Agile, Scrum, internacionalización (i18n / l10n), autenticación y autorización, MFA, cifrado en reposo, SaaS, mentoría

\cvsect{Experiencia Laboral}

\begin{entrylist}
	% Ver nota en resume.tex: el tramo Nov 2024 - May 2026 se separa a propósito, para
	% que se lea como recorrido (construir -> tracción -> constituir la sociedad) y no
	% como un periodo suelto de freelance entre empleos.
	\entry
		{May 2026 -- Actualidad}
		{Cofundador y CEO}
		{The Lovelust Company, SL}
		{Cofundé la sociedad que opera \textbf{LoveLust} y \textbf{Zierzo}, y la dirijo como CEO sin dejar de ser su responsable técnico -- LoveLust llega a los \textbf{16.000 usuarios activos mensuales} con monetización por suscripción. Asumo ambos productos \textbf{de extremo a extremo y a escala mundial}, además de \textbf{marketing}, finanzas y \textbf{cumplimiento normativo} (RGPD). Construí la plataforma que los sostiene: \textbf{microservicios NestJS} (gateway, sincronización, recompensas, notificaciones push, webhooks de RevenueCat, borrado RGPD, encuestas y soporte con pedidos de Shopify), webs en \textbf{Next.js} (web pública, back office y encuestas), PostgreSQL y un pipeline de analítica con Grafana/Prometheus. LoveLust está traducida a \textbf{14 idiomas}. Uso \textbf{Claude}, LLMs y \textbf{MCP} para acelerar el desarrollo, \textbf{automatizar pruebas} y generar assets y contenido de las apps. Represento a la empresa ante socios, prescriptores y organizaciones externas, y en eventos del sector. \\
		\texttt{React Native}\slashsep\texttt{Expo}\slashsep\texttt{TypeScript}\slashsep\texttt{RevenueCat}\slashsep\texttt{HealthKit/Health Connect}\slashsep\texttt{EAS}\slashsep\texttt{NestJS}\slashsep\texttt{Next.js}\slashsep\texttt{PostgreSQL}\slashsep\texttt{Claude}\slashsep\texttt{MCP}}
	\entry
		{Nov 2024 -- May 2026}
		{Desarrollador de Software Independiente}
		{LoveLust \slashsep Zierzo}
		{Diseñé, desarrollé y publiqué dos productos desde cero hasta \textbf{App Store y Google Play} como único desarrollador, con \textbf{cifrado de datos en local} (claves en el \textbf{Keychain de iOS} y el \textbf{Keystore de Android}). Asumí todas las capas: \textbf{arquitectura}, \textbf{CI/CD}, \textbf{tests automatizados}, revisión en tiendas y despliegues progresivos. También desarrollé \textbf{sitios web para pequeños negocios locales}. La tracción de ambos productos llevó a cofundar \textbf{The Lovelust Company, SL} en mayo de 2026. \\
		\texttt{React Native}\slashsep\texttt{Expo}\slashsep\texttt{TypeScript}\slashsep\texttt{RevenueCat}\slashsep\texttt{EAS}\slashsep\texttt{Firebase}\slashsep\texttt{i18next}}
	\entry
		{Oct 2022 -- Nov 2024}
		{Ingeniero de Software}
		{Aimsun}
		{Lideré de extremo a extremo los módulos centrales de una \textbf{plataforma SaaS} -- autenticación, cuentas, suscripciones, licencias y MFA -- dentro de un equipo multifuncional. Desarrollé herramientas internas para \textbf{facturación}, \textbf{permisos} y administración de usuarios. \\
		\texttt{React}\slashsep\texttt{TypeScript}\slashsep\texttt{Next.js}\slashsep\texttt{i18next}\slashsep\texttt{MapBox}\slashsep\texttt{GitLab CI/CD}\slashsep\texttt{Cypress}}
	\entry
		{Dic 2021 -- Oct 2022}
		{Desarrollador de Software}
		{Viewnext}
		{Desarrollé una \textbf{plataforma de reservas de puntos de carga para vehículos eléctricos}, con mapa interactivo y gestión de reservas. Contribuí a una \textbf{biblioteca de componentes compartida} con requisitos estrictos de calidad: cobertura de tests casi total y ausencia de code smells. Fui referente técnico de desarrolladores junior. \\
		\texttt{Angular}\slashsep\texttt{TypeScript}\slashsep\texttt{RxJS}\slashsep\texttt{Jenkins}\slashsep\texttt{GitLab CI/CD}}
	\entry
		{May 2021 -- Dic 2021}
		{Desarrollador de Software}
		{adidas}
		{Reemplacé herramientas internas legacy por un \textbf{editor de contenidos} moderno. Introduje la \textbf{internacionalización} y optimicé los flujos de traducción en toda la plataforma. \\
		\texttt{React}\slashsep\texttt{TypeScript}\slashsep\texttt{Redux}\slashsep\texttt{i18n}}
	\entry
		{Dic 2019 -- May 2021}
		{Desarrollador de Software}
		{Telefónica / Eleven Paths}
		{Lideré la migración de una plataforma de gobernanza y gestión de riesgos a una arquitectura moderna con \textbf{Angular}. Entregué funcionalidades empresariales y mantuve servicios backend de informes. \\
		\texttt{Angular}\slashsep\texttt{TypeScript}\slashsep\texttt{RxJS}\slashsep\texttt{Java}\slashsep\texttt{i18n}}
	\entry
		{Oct 2017 -- Dic 2019}
		{Desarrollador de Software}
		{Ingenia Sistemas Avanzados para el Transporte}
		{Desarrollé dos productos para empresas de autobuses: una \textbf{plataforma de venta de billetes personalizable por operador} con descuentos y múltiples métodos de pago, y una \textbf{herramienta de back-office interna} para gestionar estaciones, líneas, conductores y vehículos con rutas en \textbf{Google Maps}. \\
		\texttt{Angular}\slashsep\texttt{TypeScript}\slashsep\texttt{RxJS}\slashsep\texttt{Cordova}\slashsep\texttt{Stripe}\slashsep\texttt{i18n}}
	\entry
		{Mar 2017 -- Sep 2017}
		{Desarrollador de Software}
		{Rananegra}
		{Webs con CMS y una herramienta de emails automáticos con suscripciones. \\
		\texttt{Symfony}\slashsep\texttt{Angular}\slashsep\texttt{TypeScript}\slashsep\texttt{Stripe}}
\end{entrylist}

\cvsect{Proyectos}

% Las apps publicadas son la prueba más fuerte del CV: sección propia con enlaces que se
% abren en el móvil. get.* redirige a la tienda que corresponda.

\begin{entrylist}
	\entry
		{App móvil}
		{LoveLust}
		{\href{https://get.lovelust.health}{get.lovelust.health}}
		{App de salud sexual: \textbf{16.000 usuarios activos mensuales}, suscripciones, HealthKit/Health Connect y 14 idiomas.}
	\entry
		{App móvil}
		{Zierzo}
		{\href{https://get.zierzo.app}{get.zierzo.app}}
		{Autobuses, tranvía, Bizi y cines de Zaragoza en una app, en tiempo real.}
	\entry
		{Código abierto}
		{API Canopus}
		{\href{https://github.com/endworks/canopus}{github.com/endworks/canopus}}
		{\textbf{Microservicios NestJS} con API gateway, base de Zierzo: transporte urbano, cines, tiempo, diccionario de la RAE y push. Turborepo, Docker, imágenes firmadas y CI/CD por servicio.}
\end{entrylist}

\cvsect{Certificaciones}

% Separadas de Formación y colocadas primero: son las credenciales más recientes y
% relevantes, y antes aparecían las últimas, por debajo de dos titulaciones de
% 2012-2016, en una lista cuyo orden era 2014, 2012, 2025, 2025.
%
% Emisores corregidos: "Formación San Miguel" y "Kings Corner English School" son las
% academias que prepararon los exámenes, NO los organismos certificadores.

\begin{entrylist}
	\shortentry
		{2025}
		{Google Cloud Professional Cloud Architect}
		{Google Cloud}
		{}

	\shortentry
		{2025}
		{Oxford Test of English -- Advanced (C1)}
		{Oxford University Press}
		{}
\end{entrylist}

\cvsect{Formación}

% CORRECCIÓN IMPORTANTE: DAM es un Ciclo Formativo de Grado Superior -> "Técnico
% Superior"; SMR es un Ciclo Formativo de Grado Medio -> "Técnico". En España "Grado"
% designa un título universitario, así que la redacción anterior ("Grado en...")
% atribuía una titulación que no se posee -- algo que cualquier recruiter español
% detecta de inmediato. No hay ninguna desventaja en el título correcto: un CFGS en DAM
% es una vía plenamente reconocida, y con 8 años de experiencia nadie evalúa esto.

\begin{entrylist}
	\shortentry
		{2014 -- 2016}
		{Técnico Superior en Desarrollo de Aplicaciones Multiplataforma}
		{CPIFP Los Enlaces}
		{}

	\shortentry
		{2012 -- 2014}
		{Técnico en Sistemas Microinformáticos y Redes}
		{CPIFP Los Enlaces}
		{}
\end{entrylist}

% Pie a DOS columnas, no tres. La columna "Soft Skills" desaparece: su encabezado
% seguía en INGLÉS en el CV español (deriva heredada de copiar el fichero), y sus
% contenidos eran o bien no verificables o bien universales. Lo sustantivo (Agile/Scrum,
% mentoría, liderazgo) vive ahora en Competencias Técnicas > Metodologías.

% A todo el ancho desde que se añadió Proyectos: el pie a dos columnas no podía partirse
% entre páginas y empujaba el CV a una tercera página.
\cvsect{Idiomas}
\textbf{Español} y \textbf{rumano} -- nativos \quad\textbullet\quad \textbf{Inglés} -- C1 (Oxford Test of English, 2025)

\cvsect{Voluntariado}
\textbf{Fundador y Presidente}, Asociación Juvenil Brovrashgka (sin ánimo de lucro): \textbf{presupuesto}, subvenciones, \textbf{voluntarios} y eventos públicos.

\end{document}
